\begin{prob}
    Express the rate of growth of the following function using $\Theta$-notation, and
    prove your answer by finding constants which satisfy the definition
    of $\Theta$-notation. Note: please do \emph{not} use limits to prove your answer (though
    you can use them to check your work).

    \[
        f(n) = \frac{n^2 + 2n - 5}{n - 10}
    \]
    
    \begin{soln}
        We will show that \( f(n) = \Theta(n) \).

        Recall that \( f(n) = \Theta(g(n)) \) if there exist positive constants \( c_1, c_2, \) and \( N \) such that
        \[
            c_1 g(n) \leq f(n) \leq c_2 g(n) \quad \text{for all } n \geq N.
        \]
        In this case, we will prove the inequality for $g(n) = n$. There are infinitely many choices of $c_1$, $c_2$, and $N$ that will satisfy the inequality. We will demonstrate such a proof:



        \underline{Upper Bound}: \( f(n) \leq c_2 n \)

        \begin{align*}
            f(n) &= \frac{n^2 + 2n - 5}{n - 10} \\
            &\leq \frac{n^2 + 2n}{n - 10} \quad \text{(since } -5 \leq 0) \\
            &= \frac{n(n + 2)}{n - 10} \\
            \intertext{For \( n \geq 22 \), \( n - 10 \geq \dfrac{n}{2} \) (since \( n - 10 \geq 12 \) and \( n/2 = 11 \)), so:}
            &\leq \frac{n(n + 2)}{n/2} = 2(n + 2) \\
            &\leq 2n + 4 \\
            &\leq 3n \quad \text{for all } n \geq 4
        \end{align*}

        Therefore, we can choose \( c_2 = 3 \) and \( N \geq 22 \) for the upper bound.

        \vspace{1em}
        \underline{Lower Bound}: \( f(n) \geq c_1 n \)

        \begin{align*}
            f(n) &= \frac{n^2 + 2n - 5}{n - 10} \\
            &\geq \frac{n^2 + 2n - 5}{n} \quad \text{(since } n - 10 \leq n) \\
            &= n + 2 - \frac{5}{n} \\
            &\geq n + 2 - 1 \quad \text{for } n \geq 5 \quad \left( \text{since } \frac{5}{n} \leq 1 \right) \\
            &= n + 1 \\
            &\geq n \quad \text{for all } n \geq 5
        \end{align*}

        Therefore, we can choose \( c_1 = 1 \) and \( N \geq 5 \) for the lower bound.

        \vspace{1em}
        \underline{Conclusion}:

        Combining both bounds, we have for \( n \geq 22 \):
        \[
            n \leq f(n) \leq 3n
        \]
        Hence, \( f(n) = \Theta(n) \) with constants \( c_1 = 1 \), \( c_2 = 3 \), and \( N = 22 \).
    \end{soln}
    
\end{prob}
